% !TeX program = xelatex                                      % 使用 XeLaTeX 编译中文报告。
\documentclass[UTF8,a4paper,11pt]{ctexart}                    % 中文文章类、A4 纸、11 磅字号。
\usepackage[left=2cm,right=2cm,top=2cm,bottom=2cm]{geometry}  % 设置紧凑页边距。
\usepackage{amsmath,amssymb,bm}                               % 数学公式、符号和向量粗体。
\usepackage{graphicx}                                         % 插入 Python 生成的图片。
\usepackage{booktabs}                                         % 绘制三线表。
\usepackage{float}                                            % 使用 [H] 固定图表位置。
\usepackage{fancyhdr}                                         % 设置页眉。
\usepackage{xcolor}                                           % 设置 Python 代码颜色。
\usepackage{listings}                                         % 在附录中排版 Python 代码。
\usepackage[hidelinks]{hyperref}                              % 提供引用且隐藏链接边框。

\geometry{headheight=14pt}                                    % 为页眉预留足够高度。
\setlength{\parindent}{2em}                                   % 中文段落首行缩进。
\setlength{\parskip}{0.2em}                                   % 设置较小段间距。
\pagestyle{fancy}                                             % 启用页眉。
\fancyhf{}                                                    % 清除默认页眉页脚。
\fancyhead[L]{数值分析期末上机报告\quad 学号：241098117\quad 姓名：张城宗} % 固定个人信息。
\fancyhead[R]{\thepage}                                       % 页眉右侧显示页码。

\lstdefinestyle{python}{                                      % 定义 Python 代码样式。
  language=Python,                                            % 启用 Python 语法高亮。
  basicstyle=\ttfamily\footnotesize,                          % 使用紧凑等宽字体。
  keywordstyle=\color{blue!75!black},                         % 关键字显示为蓝色。
  commentstyle=\color{green!45!black},                        % 注释显示为绿色。
  stringstyle=\color{orange!70!black},                        % 字符串显示为橙色。
  numbers=left,                                               % 左侧显示行号。
  numberstyle=\tiny\color{gray},                              % 行号使用灰色小字。
  frame=single,                                               % 为代码添加单线边框。
  breaklines=true,                                            % 长代码自动换行。
  showstringspaces=false                                      % 不显示字符串空格标记。
}

\begin{document}                                              % 报告正文开始。

\begin{center}                                                % 标题居中。
  {\LARGE\bfseries 数值分析期末上机报告}

  \vspace{0.3em}
  学号：241098117\qquad 姓名：张城宗
\end{center}

% ======================================================================
\section{题目一：单步法与 Adams PECE 方法比较}

\subsection{问题描述}

求解初值问题
\[
  y'=-\frac{1}{x^2}-\frac{y}{x}-y^2,\qquad
  1\le x\le 2,\qquad y(1)=-1,
\]
取固定步长 $h=0.1$。

\begin{enumerate}
  \item 分别使用向前 Euler、改进 Euler、Heun、中点和经典四阶 Runge--Kutta
        方法计算数值解，并比较不同方法的结果。
  \item 使用 RK4 提供起步值，实现四阶 Adams PECE 预测校正方法，并与 RK4 比较。
\end{enumerate}

\subsection{方法介绍}

向前 Euler 方法采用当前节点斜率推进：
\[
  y_{n+1}=y_n+h f(x_n,y_n),
\]
它实现简单，但只有一阶精度。

改进 Euler、Heun 和中点方法均属于二阶显式 Runge--Kutta 方法。以改进 Euler 为例，
\[
  K_1=f(x_n,y_n),\qquad K_2=f(x_n+h,y_n+hK_1),\qquad
  y_{n+1}=y_n+\frac{h}{2}(K_1+K_2).
\]

经典 RK4 方法为
\[
  \begin{aligned}
    K_1&=f(x_n,y_n),&
    K_2&=f(x_n+h/2,y_n+hK_1/2),\\
    K_3&=f(x_n+h/2,y_n+hK_2/2),&
    K_4&=f(x_n+h,y_n+hK_3),\\
    y_{n+1}&=y_n+\frac{h}{6}(K_1+2K_2+2K_3+K_4).
  \end{aligned}
\]
它每步需要四次函数计算，理论收敛阶为四阶。

四阶 Adams PECE 方法先用四步 Adams--Bashforth 公式预报，
\[
  y_{n+1}^{P}=y_n+\frac{h}{24}
  (55f_n-59f_{n-1}+37f_{n-2}-9f_{n-3}),
\]
再用 Adams--Moulton 公式校正，
\[
  y_{n+1}=y_n+\frac{h}{24}
  \left(9f(x_{n+1},y_{n+1}^{P})+19f_n-5f_{n-1}+f_{n-2}\right).
\]
由于多步法需要历史值，前三步使用 RK4 起步。

\subsection{数值结果}

\begin{table}[H]
  \centering
  \caption{不同方法在终点 $x=2$ 处的数值解}
  \begin{tabular}{lc}
    \toprule
    方法 & 终点数值解 $y(2)$ \\
    \midrule
    向前 Euler & $-3.5510956192$ \\
    改进 Euler & $-5.1032101606$ \\
    Heun & $-5.0041974679$ \\
    中点 & $-4.9557885785$ \\
    RK4 & $-5.4009915651$ \\
    Adams PECE & $-5.3802222860$ \\
    \bottomrule
  \end{tabular}
\end{table}

Adams PECE 与 RK4 在终点的绝对差为
\[
  |y_{\mathrm{PECE}}(2)-y_{\mathrm{RK4}}(2)|
  \approx 2.077\times 10^{-2}.
\]

\IfFileExists{python_p1_solutions.pdf}{
  \begin{figure}[H]
    \centering
    \includegraphics[width=0.78\textwidth]{python_p1_solutions.pdf}
    \caption{五种单步法与 Adams PECE 数值解比较}
  \end{figure}
}{}

\subsection{结果分析}

\begin{enumerate}
  \item 本题解在区间后半段下降较快，因此不同格式在步长 $h=0.1$ 下产生了较明显差异。
        向前 Euler 只有一阶精度，与高阶方法结果差异最大。
  \item 三种二阶方法的结果较为接近，但仍与 RK4 存在明显差异，说明该步长对于本题并不算小。
  \item Adams PECE 使用 RK4 起步，其结果与 RK4 最接近。它在主循环中复用历史函数值，
        每步函数计算次数少于 RK4，适合固定步长下的连续计算。
  \item 原数值 ODE 报告中本题终点结果约为 $-0.6$，但该结果与报告所写方程不一致；
        本报告中的数据由附录 Python 程序按上述方程重新计算得到。
\end{enumerate}

% ======================================================================
\clearpage
\section{题目二：二阶方程化为一阶系统并验证 RK4 收敛阶}

\subsection{问题描述}

求解二阶常微分方程初值问题
\[
  y''+\frac{2}{x}y'-\frac{6}{x^2}y=5-6x+7x^2,\qquad 1<x<2,
\]
\[
  y(1)=\frac12,\qquad y'(1)=2.
\]
精确解为
\[
  y(x)=x^2-x^3+\frac12x^4+x^2\ln x.
\]
分别取步长
\[
  h=\frac1{200},\frac1{100},\frac1{50},\frac1{25},
  \frac1{10},\frac15,\frac14,
\]
使用经典 RK4 方法计算最大绝对误差并分析收敛阶。

\subsection{方法介绍}

令
\[
  w_1=y,\qquad w_2=y',
\]
则二阶方程可化为一阶系统
\[
  \bm w'=
  \begin{bmatrix}
    w_2\\
    -\dfrac{2}{x}w_2+\dfrac{6}{x^2}w_1+5-6x+7x^2
  \end{bmatrix},
  \qquad
  \bm w(1)=
  \begin{bmatrix}
    1/2\\2
  \end{bmatrix}.
\]
对向量状态 $\bm w$ 使用 RK4 时，计算公式与标量情形相同。第一分量 $w_1$ 即为原方程的数值解。

定义最大绝对误差
\[
  E(h)=\max_n|w_{1,n}-y(x_n)|.
\]
对相邻两个步长，数值收敛阶估计为
\[
  p\approx
  \frac{\log(E(h_1)/E(h_2))}{\log(h_1/h_2)}.
\]

\subsection{数值结果}

\begin{table}[H]
  \centering
  \caption{不同步长下 RK4 最大误差与估计收敛阶}
  \begin{tabular}{ccc}
    \toprule
    步长 $h$ & 最大绝对误差 $E(h)$ & 估计收敛阶 $p$ \\
    \midrule
    $1/200$ & $2.4661\times10^{-10}$ & --- \\
    $1/100$ & $3.9227\times10^{-9}$  & $3.9915$ \\
    $1/50$  & $6.1913\times10^{-8}$  & $3.9803$ \\
    $1/25$  & $9.6405\times10^{-7}$  & $3.9608$ \\
    $1/10$  & $3.4751\times10^{-5}$  & $3.9123$ \\
    $1/5$   & $4.8875\times10^{-4}$  & $3.8140$ \\
    $1/4$   & $1.1220\times10^{-3}$  & $3.7243$ \\
    \bottomrule
  \end{tabular}
\end{table}

\IfFileExists{python_p2_solution.pdf}{
  \begin{figure}[H]
    \centering
    \includegraphics[width=0.48\textwidth]{python_p2_solution.pdf}
    \includegraphics[width=0.48\textwidth]{python_p2_convergence.pdf}
    \caption{RK4 数值解与误差收敛阶}
  \end{figure}
}{}

\subsection{结果分析}

\begin{enumerate}
  \item 当步长由 $1/4$ 缩小至 $1/200$ 时，最大误差由约 $10^{-3}$ 降至约 $10^{-10}$，
        表明 RK4 对该光滑问题具有较高精度。
  \item 在较小步长下，估计收敛阶接近 $4$，与经典 RK4 的理论四阶精度一致。
        较大步长下估计阶略低，是因为渐近误差关系 $E(h)\approx Ch^4$ 尚未完全成立。
  \item 将高阶方程化为一阶系统后，可以直接复用向量 RK4 程序；
        这种状态变量方法也适用于更高阶方程和一般常微分方程组。
\end{enumerate}

\subsection{结论}

两道题均使用 Python 完成数值计算。第一题说明不同单步法和多步法在精度与计算量上存在差异；
第二题验证了经典 RK4 对光滑问题具有四阶收敛性。实际计算时应根据精度要求、稳定性以及函数计算成本选择方法和步长。

% ======================================================================
\clearpage
\section*{附录：Python 程序}
\addcontentsline{toc}{section}{附录：Python 程序}

\subsection*{A.1\quad 题目一 Python 程序}

\begin{lstlisting}[style=python,caption={题目一：五种单步法与 Adams PECE 比较}]
import numpy as np  # 导入 NumPy；用于创建网格、预分配数组和进行向量运算。
import matplotlib.pyplot as plt  # 导入 Matplotlib；用于绘制数值解比较图。

f = lambda x, y: -1 / x**2 - y / x - y**2  # 定义 ODE 右端 f(x,y)；** 表示乘方。
x0, x_end, y0, h = 1.0, 2.0, -1.0, 0.1  # 设置区间、初值和固定步长。
N = round((x_end - x0) / h)  # 计算固定步长下的总推进步数。
x = x0 + np.arange(N + 1) * h  # 构造等距节点 x_0,...,x_N。

def rk4(f, x, y0, h):  # 定义经典四阶 Runge--Kutta 求解函数。
    y = np.zeros(x.size)  # 预分配数值解数组，避免循环中反复扩容。
    y[0] = y0  # 写入初值 y(x_0)=y_0。
    for n in range(x.size - 1):  # 从左端点逐步推进到右端点。
        k1 = f(x[n], y[n])  # 第一阶段使用当前点斜率。
        k2 = f(x[n] + h/2, y[n] + h*k1/2)  # 第二阶段在预测中点计算斜率。
        k3 = f(x[n] + h/2, y[n] + h*k2/2)  # 第三阶段再次修正中点斜率。
        k4 = f(x[n] + h, y[n] + h*k3)  # 第四阶段在右端预测点计算斜率。
        y[n+1] = y[n] + h*(k1 + 2*k2 + 2*k3 + k4)/6  # 加权组合四个阶段，得到四阶近似。
    return y  # 返回全部节点上的 RK4 数值解。

y_euler = np.zeros(N + 1)  # 预分配向前 Euler 数值解。
y_improved = np.zeros(N + 1)  # 预分配改进 Euler 数值解。
y_heun = np.zeros(N + 1)  # 预分配二阶 Heun 数值解。
y_midpoint = np.zeros(N + 1)  # 预分配中点方法数值解。
for y in (y_euler, y_improved, y_heun, y_midpoint):  # 依次访问四个数值解数组。
    y[0] = y0  # 为每种单步法写入相同初值。

for n in range(N):  # 对四种单步法同时逐步推进。
    y_euler[n+1] = y_euler[n] + h*f(x[n], y_euler[n])  # 向前 Euler 是一阶显式方法。
    k1 = f(x[n], y_improved[n])  # 改进 Euler 的左端斜率。
    k2 = f(x[n] + h, y_improved[n] + h*k1)  # 在 Euler 预测点计算右端斜率。
    y_improved[n+1] = y_improved[n] + h*(k1+k2)/2  # 两端斜率平均得到二阶结果。
    k1 = f(x[n], y_heun[n])  # Heun 方法的第一阶段斜率。
    k2 = f(x[n] + 2*h/3, y_heun[n] + 2*h*k1/3)  # 在 2/3 步位置计算第二阶段。
    y_heun[n+1] = y_heun[n] + h*(k1+3*k2)/4  # 使用满足二阶条件的权重组合。
    k1 = f(x[n], y_midpoint[n])  # 中点方法先计算当前点斜率。
    k2 = f(x[n] + h/2, y_midpoint[n] + h*k1/2)  # 用 k1 预测中点并计算中点斜率。
    y_midpoint[n+1] = y_midpoint[n] + h*k2  # 使用中点斜率推进整步。

y_rk4 = rk4(f, x, y0, h)  # 调用 RK4 函数计算四阶参考数值解。
y_pece = np.zeros(N + 1)  # 预分配 Adams PECE 数值解。
f_values = np.zeros(N + 1)  # 保存历史函数值，以供多步公式重复使用。
y_pece[:4] = y_rk4[:4]  # 使用 RK4 提供 y_0,y_1,y_2,y_3 四个起步值。
for i in range(4):  # 计算四个起步节点对应的函数值。
    f_values[i] = f(x[i], y_pece[i])  # 保存 f_i=f(x_i,y_i)。
for n in range(3, N):  # 从已有四个历史点的位置开始执行 PECE。
    y_pred = y_pece[n] + h*(55*f_values[n]-59*f_values[n-1]+37*f_values[n-2]-9*f_values[n-3])/24  # AB4 预报。
    f_pred = f(x[n+1], y_pred)  # 在预测值处计算新节点函数值。
    y_pece[n+1] = y_pece[n] + h*(9*f_pred+19*f_values[n]-5*f_values[n-1]+f_values[n-2])/24  # AM4 校正。
    f_values[n+1] = f(x[n+1], y_pece[n+1])  # 保存校正值对应的函数值，供下一步使用。

print("method          y(2)")  # 打印结果表表头。
for name, y in [("Euler", y_euler), ("Improved Euler", y_improved), ("Heun", y_heun), ("Midpoint", y_midpoint), ("RK4", y_rk4), ("Adams PECE", y_pece)]:  # 配对方法名和结果数组。
    print(f"{name:15s} {y[-1]:.10f}")  # 输出每种方法的终点数值解。
print("PECE-RK4 difference =", abs(y_pece[-1] - y_rk4[-1]))  # 输出 PECE 与 RK4 的终点差。

for name, y in [("Euler", y_euler), ("Improved Euler", y_improved), ("Heun", y_heun), ("Midpoint", y_midpoint), ("RK4", y_rk4), ("Adams PECE", y_pece)]:  # 遍历全部方法。
    plt.plot(x, y, marker="o", label=name)  # 在同一坐标系绘制各方法数值解。
plt.xlabel("x")  # 设置横轴名称。
plt.ylabel("y")  # 设置纵轴名称。
plt.grid(True)  # 打开网格，便于比较曲线。
plt.legend()  # 显示方法图例。
plt.tight_layout()  # 自动调整边距，避免文字被裁切。
plt.savefig("python_p1_solutions.pdf")  # 保存矢量 PDF，供 LaTeX 自动插入。
plt.close()  # 关闭当前图，释放绘图资源。
\end{lstlisting}

\subsection*{A.2\quad 题目二 Python 程序}

\begin{lstlisting}[style=python,caption={题目二：二阶 ODE 化为一阶系统并验证 RK4 收敛阶}]
import numpy as np  # 导入 NumPy；用于向量状态、网格、误差和对数运算。
import matplotlib.pyplot as plt  # 导入 Matplotlib；用于绘制解曲线和双对数误差图。

def F(x, w):  # 定义一阶系统右端；w[0]=y，w[1]=y'。
    return np.array([w[1], -2*w[1]/x + 6*w[0]/x**2 + 5 - 6*x + 7*x**2])  # 返回列状态的两个导数分量。

def exact(x):  # 定义题目给出的精确解。
    return x**2 - x**3 + 0.5*x**4 + x**2*np.log(x)  # NumPy 运算自动逐元素作用于数组。

def solve_rk4(h):  # 定义给定步长下的向量 RK4 求解器。
    x0, x_end = 1.0, 2.0  # 设置积分区间。
    N = round((x_end - x0) / h)  # 计算总步数。
    x = x0 + np.arange(N + 1) * h  # 创建等距节点。
    w = np.zeros((N + 1, 2))  # 每一行保存一个节点的 [y,y']。
    w[0] = [0.5, 2.0]  # 写入初始状态 [y(1),y'(1)]。
    for n in range(N):  # 对向量状态逐步执行 RK4。
        k1 = F(x[n], w[n])  # 第一阶段斜率向量。
        k2 = F(x[n] + h/2, w[n] + h*k1/2)  # 第二阶段中点斜率向量。
        k3 = F(x[n] + h/2, w[n] + h*k2/2)  # 第三阶段修正的中点斜率向量。
        k4 = F(x[n] + h, w[n] + h*k3)  # 第四阶段右端斜率向量。
        w[n+1] = w[n] + h*(k1 + 2*k2 + 2*k3 + k4)/6  # 使用 RK4 权重更新全部状态分量。
    return x, w  # 返回节点与完整状态数组。

h_values = np.array([1/200, 1/100, 1/50, 1/25, 1/10, 1/5, 1/4])  # 设置题目要求的步长。
errors = np.zeros(h_values.size)  # 预分配每个步长对应的最大绝对误差。
for i, h in enumerate(h_values):  # enumerate 同时给出数组下标和步长。
    x, w = solve_rk4(h)  # 计算当前步长下的向量数值解。
    errors[i] = np.max(np.abs(w[:, 0] - exact(x)))  # 第一分量是 y；对全部节点误差取最大值。

orders = np.full(h_values.size, np.nan)  # 第一组数据没有前一误差，因此初始填为 nan。
for i in range(1, h_values.size):  # 从第二个步长开始估计相邻误差的收敛阶。
    orders[i] = np.log(errors[i-1]/errors[i]) / np.log(h_values[i-1]/h_values[i])  # 使用一般步长比公式。

print("h          max_error          order")  # 打印误差与阶数表头。
for h, error, order in zip(h_values, errors, orders):  # zip 将三组对应数据逐行配对。
    print(f"{h:.6f}   {error:.10e}   {order:.6f}")  # 格式化输出步长、最大误差与估计阶。

x, w = solve_rk4(1/200)  # 使用最小步长生成数值解与精确解对比图。
plt.plot(x, exact(x), "k-", label="exact")  # 绘制精确解实线。
plt.plot(x, w[:, 0], "r--", label="RK4")  # 绘制 RK4 第一分量虚线。
plt.xlabel("x")  # 设置横轴名称。
plt.ylabel("y")  # 设置纵轴名称。
plt.grid(True)  # 打开网格。
plt.legend()  # 显示精确解与数值解图例。
plt.tight_layout()  # 自动调整图像边距。
plt.savefig("python_p2_solution.pdf")  # 保存解曲线图供 LaTeX 插入。
plt.close()  # 关闭解曲线图。

plt.loglog(h_values, errors, "bo-", label="RK4 error")  # 在双对数坐标绘制误差曲线。
reference = errors[0] * (h_values / h_values[0])**4  # 构造经过首个误差点的四阶参考线。
plt.loglog(h_values, reference, "k--", label="O(h^4)")  # 绘制理论四阶误差参考线。
plt.xlabel("step size h")  # 设置横轴为步长。
plt.ylabel("maximum error")  # 设置纵轴为最大绝对误差。
plt.grid(True, which="both")  # 对主刻度和次刻度均显示网格。
plt.legend()  # 显示误差曲线和参考线图例。
plt.tight_layout()  # 自动调整图像边距。
plt.savefig("python_p2_convergence.pdf")  # 保存收敛阶图供 LaTeX 插入。
plt.close()  # 关闭误差图并释放资源。
\end{lstlisting}

\end{document}                                                % 报告结束。
